\documentclass[11pt]{article}

\usepackage[utf8]{inputenc}
\usepackage[margin=1in]{geometry}
\usepackage{setspace}

\title{Plain Language Summary}

\author{Bruce Stephenson \and Robin Macomber}

\date{March 2026}

\begin{document}

\maketitle
\doublespacing

\noindent\textbf{Plain language title:} Scientists in different fields independently invented the same failure-prediction math

\bigskip

When a complex system---a heart, a power grid, a financial market, a climate system---approaches a catastrophic transition, its behavior changes in predictable ways. Small disturbances take longer to die out. Distant parts of the system begin to move in sync. Fluctuations grow. Scientists have developed mathematical tools to detect these warning signs and predict when a system is nearing a tipping point.

What is remarkable is that researchers in at least six different fields invented essentially the same mathematics for this purpose, over nine decades, mostly without knowing about each other's work. A physicist measuring ``correlation length,'' a cardiologist computing a ``DFA scaling exponent,'' a financial analyst calculating a ``Hurst exponent,'' and a machine learning engineer tuning a ``spectral radius'' are all answering the same question: how close is this system to a critical transition? Their tools detect the same signatures---growing correlations, slowing recovery, increasing sensitivity---using different notation developed independently within each discipline.

This paper surveys between six and twelve fields (depending on how strictly one counts) where this convergent discovery occurred. We classify each case as independent invention, adaptation from another field, or an early empirical observation that predated the formal mathematics. We also analyze citation networks---who cited whom---and find that cross-field awareness was remarkably low during the key period from 1987 to 2010, even though the techniques were functionally equivalent.

The convergence has practical implications. Insights from one domain can transfer to another: a power grid engineer can learn from how cardiologists detect cardiac risk, and climate scientists' tipping-point models share mathematical structure with financial crash prediction. Documenting this pattern is a step toward breaking down the disciplinary barriers that kept these communities working in parallel for decades.

\end{document}
